%============================================================
\chapter{検証障害の分類：四分類から六分類へ}
\label{part:obstacles}
%============================================================

\section{本章の目的}

第\ref{part:results}部において、族2 と族5 の双方について
一次データにあたったが、いずれも検証が完結しなかった。
族2 では個社の未使用残高が、族5 では個社の返戻金制度が必要であり、
どちらも集計統計には現れない。
本章はその取得経路を確認した結果を記録する。

結論から述べる。障害の性質は族2 と族5 で異なる。
族2 では必要な粒度のデータに到達できない。
族5 ではデータ源が想定より良好であり、307社について検定を実施できたが、
説明変数が理論の想定する範囲で変動していなかった。
\textbf{いずれもデータの追加のみでは解決しない。}

以下ではまず障害を四つに分類し、族ごとに割り当てる。
その過程で四分類に収まらない障害が二つ現れたため、末尾で追加する。

\section{検証障害の四分類}
\label{sec:obstacle-taxonomy}

「検証できない」には性質の異なる四つの状態が含まれる。
対処法が異なるため、分けて扱う必要がある。この分類は本稿の整理である。

\begin{center}
\begin{tabularx}{\textwidth}{@{}llX@{}}
\toprule
種別 & 状態 & 正しい対処 \\
\midrule
（i）測定の不在 & 世界の誰もその量を測っていない & 自ら原調査を行う。二次データでは永久に埋まらない \\
（ii）仮説の不良設定 & 理論が符号を一意に決めない & 理論を精密化する。データを増やしても解決しない \\
（iii）識別不能 & データはあるが交絡を切れない & 設計の工夫。外生変動か対照群を探す \\
（iv）アクセス制約 & 存在するが取得できない & 経路の開拓、または法的・技術的に断念 \\
\bottomrule
\end{tabularx}
\captionof{table}{検証障害の四分類。状態と正しい対処。}
\label{tab:obstacle-taxonomy}
\end{center}

（i）と（ii）は、いくらデータを集めても解決しない点で共通する。
一方（iii）と（iv）は労力や工夫で動きうる。
「検証困難」という言葉はこの四つを混ぜてしまうため、
どこで詰まっているのかを見失わせる。

\section{族2への割当}

\subsection{達成された部分}

第\ref{part:results}部で述べたとおり、$\kappa$ の水準差は媒体別パネルで測定でき、
対抗仮説（市場縮小による機械的上昇）も部分的に分離できた。
これは（iii）識別不能を、成長セグメントの存在という外生変動によって解いた例である。

\subsection{残った障害の性質}

本領域で $\phi_{\mathrm{cog}}$ が測れないのは、退蔵率が集計されていないためである。
協会の統計に対しては発行者自身から
「未回収率（期限切れ等）の状況も開示してほしい」という要望が出ており、
この統計の枠内では測定が行われていない。

これは（i）測定の不在にみえる。しかし障害の所在はより限定的である。
命題\ref{prop:phicog-requirements}の三条件のうち、
本領域で欠けているのは (2) 個票の実現利用と (3) 注意を強制する外生イベントであり、
(1) 契約上の権利は媒体別の使用期限として観測できる。
すなわち\textbf{測定手法は確立しており、欠けているのはデータの粒度である}。

したがって第\ref{sec:obstacle-taxonomy}節の分類では、
族2の $\phi_{\mathrm{cog}}$ は（iv）アクセス制約に属する。
ただし通常のアクセス制約と異なり、必要な粒度のデータは
発行者の内部にしか存在せず、公表を待つ経路がない。

個社の未使用残高が守秘義務により非公開である点は副次的である。
仮に個社データが得られても、報告項目に個票の利用記録が含まれない以上
$\phi_{\mathrm{cog}}$ は推定できない。

\begin{center}
\begin{tabularx}{\textwidth}{@{}lX@{}}
\toprule
族2 の主要障害 & （iv）アクセス制約。個票と外生ショックを欠く \\
\bottomrule
\end{tabularx}
\captionof{table}{族2の主要障害の判定。}
\label{tab:family2-obstacle}
\end{center}

\section{族5への割当}

\subsection{データ源は想定より良好であった}
\label{sec:jinzai-structure}

人材サービス総合サイトの事業所詳細ページを確認した結果、
第\ref{part:results}部で懸念した非構造化の問題は、返戻金の内容を除いて生じていなかった。

\begin{center}
\begin{tabularx}{\textwidth}{@{}lXl@{}}
\toprule
変数 & 掲載形式 & 判定 \\
\midrule
就職者数（無期／有期／日雇の別） & 数値、令和2〜7年度の6年分 & 構造化 \\
6か月以内離職者数（無期雇用） & 数値、同6年分 & 構造化 \\
離職が判明せず（無期雇用） & 数値、同6年分 & 構造化 \\
職種別の手数料実績率 & パーセント、令和7年度から & 構造化 \\
返戻金制度の有無 & 有／無 & 構造化 \\
返戻金の率・期間 & 事業者が任意形式でPDF添付 & 非構造化 \\
\bottomrule
\end{tabularx}
\captionof{table}{人材サービス総合サイトの掲載変数と構造化の状況。}
\label{tab:jinzai-variables}
\end{center}

さらに\textbf{詳細検索により、職種・手数料実績率・離職率を条件として指定でき、
検索結果に三者が同時に表示される}。
これにより事業所を一件ずつ開く必要がなくなり、
第\ref{sec:predX-result}節の307社の取得が可能になった。
ただし条件に指定できる職種は医療・介護・保育分野の9つに限られる。

三点、当初の想定を上回った。

第一に、\textbf{常用と臨時日雇が分離されている}。
第\ref{part:results}部で医師・看護師の集計値が使えなかった混入の問題は、
個社レベルでは生じない。

第二に、\textbf{手数料が率（パーセント）で公開されている}。
第\ref{part:results}部で「年収の分母がないため識別力がない」と述べた問題は、
すでに割り算された形で掲載されているため、外部から年収データを調達する必要がない。
一例として、ある事業所では令和7年度の実績率が
ソフトウェア開発40.1\%、デザイナー37.4\%、営業37.7\%と職種別に掲載されていた。

第三に、\textbf{「離職が判明せず」という欄が存在する}。
これは事業者が就職者のその後を観測できているかを示す数値であり、
第\ref{sec:info}章の $\mathcal{G}$（検証可能な共有情報）の広さの、
代理変数ではなくほぼ直接の観測である。
第\ref{part:results}部で「B1側の代理変数を作れていないのが設計上の最大の弱点」と
記した問題への解答になりうる。ただし自己申告であり、
追跡していない事業者がゼロと記入している可能性は排除できない。

また、適正認定事業者の個社ページには
「厚生労働省 人材サービス総合サイトURL/許可番号」欄があり、
認定一覧から事業所詳細ページへの連結キーが存在することを確認した。

\subsection{説明変数が制度により規定される範囲}
\label{sec:regulated-regressor}

検定に用いようとする説明変数は\textbf{返戻金制度の強度 $b$}、
具体的には返戻率と対象期間である。
このうち\emph{対象期間}が、一部の分野で制度により下限を課されている。

医療・介護・保育分野の適正認定制度では、2024年度認定より、
早期離職時の返戻金は求職者が就職後少なくとも6か月以内に離職した場合を
対象としていることが認定条件となっている。

\textbf{この分野の認定事業者に限り、対象期間の下限は市場選択ではなく規制の産物である。}
認定事業者を標本に含めたまま対象期間を説明変数とすると、
測っているのは規制の遵守であって契約選択ではない。
ただしこの論点が及ぶ範囲は限定的である（注\ref{rem:regulated-scope}）。

\begin{remark}[適用範囲の限定]
\label{rem:regulated-scope}
この論点が及ぶ範囲は限定的である。

第一に、対象は医療・介護・保育分野の\emph{認定}事業者のみである。
認定事業者は実企業で60--70程度、
事業所30,561のうちごく一部にすぎない。

第二に、規定されるのは対象期間の\emph{下限}のみであり、
返戻率の水準や下限を超える期間の設定は規制されない。

第三に、認定は任意であり、認定を受けない選択も可能である。
ただし2026年8月に配置基準の柔軟化等に係る特例対応の要件として
適正認定事業者の活用が含まれたため、認定に経済的価値が付与された。
遵守圧力は強まる方向にある。

したがって「制度が説明変数を規定する」という主張は、
\textbf{標本の一部について、変数の一成分の下限にのみ当てはまる}。
これを分野全体や返戻金制度一般に拡張することはできない。
\end{remark}

これは（iii）識別不能に属するが、同時に識別の機会でもある。
規制が特定分野にのみかかっているため、三群の比較が構成できる。

\begin{center}
\begin{tabular}{@{}llll@{}}
\toprule
群 & 分野 & 認定 & 返戻金の決定要因 \\
\midrule
A & 医療・介護・保育 & あり & 規制（6か月が下限） \\
B & 医療・介護・保育 & なし & 市場選択 \\
C & その他分野 & --- & 市場選択 \\
\bottomrule
\end{tabular}
\captionof{table}{返戻金制度の決定要因による三群（A/B/C）の比較。}
\label{tab:predB-groups}
\end{center}

B対C が本来の仮説検定にあたる。
A対B は別の問いに答える。すなわち\textbf{規制が市場をどれだけ動かしたか}である。
B群も大半が6か月を提供しているなら規制は非拘束的で、市場が既に選んでいたことの追認にすぎない。
B群が1〜3か月中心なら規制は拘束的で、市場均衡から離れた水準を強制したことになる。

A対B は当初の仮説とは別物だが、
制度層が $\Phi$ をどれだけ書き換えられるかの直接測定であり、
第\ref{sec:layers}章の主張の検証になる。

\subsection{記録の粒度が二つの資料で異なる}
\label{sec:unit-mismatch}

より実務的だが深刻な問題がある。
返戻金制度と手数料率は別の資料から取得する必要があるが、
両者は記録の粒度が異なる。

\begin{center}
\begin{tabularx}{\textwidth}{@{}lXl@{}}
\toprule
資料 & 内容 & 記録の粒度 \\
\midrule
適正認定事業者一覧 & 認定を受けた事業者の名簿。厚生労働省が公表 & 企業 $\times$ 分野 \\
人材サービス総合サイト & 就職者数、離職者数、職種別手数料率 & 事業所 \\
\bottomrule
\end{tabularx}
\captionof{table}{返戻金制度と手数料率、それぞれの資料と記録粒度。}
\label{tab:record-granularity}
\end{center}

認定事業者一覧では、同一企業が医療・介護・保育の三分野で別レコードとして登録され、
認定番号も分野ごとに異なる。
一方、人材サービス総合サイトの許可番号は特定の事業所を指す。

したがって、複数事業所を持つ企業では、認定に紐づく許可番号が
企業全体の実績をカバーしているかが不明である。
ある大手は介護分野で9職種、医療分野で6職種の認定を受けているが、
事業所単位の職種別手数料率がその粒度をカバーする保証はない
（第\ref{sec:jinzai-structure}節で確認した事業所は5職種のみであった）。

\textbf{この誤差は規模と相関する}。大手ほど不整合が大きく、
認定事業者には大手が多い。系統的なバイアスになる。

\subsection{職種分類の三系統分裂}

接続を要する分類が三つあり、対応表が存在しない。

\begin{center}
\begin{tabularx}{\textwidth}{@{}llX@{}}
\toprule
系統 & 区分 & 例 \\
\midrule
職業紹介事業報告書 & 約110職種 & 看護師、准看護師 \\
人材サービス総合サイト & コード番号制 & 009、010、017、033、048 \\
適正認定制度 & 分野別の細区分 & 看護職、リハビリテーション専門職 \\
\bottomrule
\end{tabularx}
\captionof{table}{職種分類の三系統と、それぞれの区分方式。}
\label{tab:occupation-classification-systems}
\end{center}

認定制度側は「看護職」が看護師・准看護師・保健師・助産師を含むと注記しているが、
他は不明である。コード番号制の体系は職業紹介事業報告書の職種順序と一致しない。
対応表の作成自体は労力で解決可能な範囲にある。

\subsection{認定事業者に限った分析の制約}
\label{sec:nintei-small}

適正認定事業者一覧は約100件である。
認定は企業$\times$分野単位であり同一企業が複数分野で重複するため、
実企業数は60から70程度と見込まれる。

\textbf{認定の有無を説明変数とする分析は、この規模では成立しない。}
第\ref{sec:predX-result}節の標本307社のうち認定事業者は一部にすぎず、
群間比較の検定力が得られない。
第\ref{sec:frame}節で述べたとおり、フレームの網羅性は標本規模を保証しない。

ただしこれは\textbf{認定を変数とする分析に限った制約}である。
手数料実績率と離職率の関係は、認定の有無によらず全事業者について測定でき、
実際に307社で実施した。

\subsection{可測性仮説は符号が一意に定まらない}
\label{sec:predB-illposed}

族5 で検証しようとするのは、第\ref{sec:info}章の可測性が契約形態を決めるという主張である。
素朴に述べれば「成果の事後検証が難しい職種ほど返戻金制度が強い」となるが、
この形では逆向きの力を無視している。

例\ref{ex:bstar}の最適強度 $b^\star \propto 1/(1+r\sigma^2 k)$ から二つの力が出る。

\begin{itemize}
\item \textbf{B1（分子・インセンティブ側）}：事業者の努力が定着に効きやすい職種ほど、
$\pi$ を $y$ に載せる価値が高い。返戻金は\emph{強く}なる。
\item \textbf{B2（分母・リスク側）}：離職が事業者の制御外の要因で決まる職種ほど、
$y$ に載せてもリスクを移転するだけである。返戻金は\emph{弱く}なる。
\end{itemize}

素朴な定式化はB1のみを見ている。多くの職種で両者は同方向に共変するため、
\textbf{符号が一意に決まらない}。

\begin{remark}[これは分野全体の未解決問題である]
符号が定まらないことは、本稿に固有の設定上の不備ではない。

\textbf{リスク・インセンティブ・トレードオフ（RIT）は四半世紀の論争の対象であり、
実証的には支持されていない。} \cite{prendergast1999,prendergast2002} は
26の実証研究を検討し、理論が予測する負の関係を見出したのは4件のみであり、
むしろ不確実性と成果連動の間に\emph{正}の相関が観察されると報告している。
Prendergast 自身の説明は、不確実な状況では企業は労働者の時間の使い方を把握できないため
権限を委譲し、その裁量を制約するために産出に基づく報酬を用いる、というものである。
すなわち本節でB1と呼んだ力は、既に定式化された対抗理論として存在する。

論争は現在も継続している。実験室では支持する結果が報告されはじめており、
期待効用理論を離れて Rank-Dependent Utility を用いると RIT が頑健に成立するとの
主張もある。

したがってこの仮説は第\ref{sec:obstacle-taxonomy}節の分類における
（ii）仮説の不良設定に属する。
検定可能な形にするには、B1とB2を分離して操作化しなければならない。
\end{remark}

\subsection{価格側からの検定}
\label{sec:what-to-measure}

$b^\star$ の式そのものは、必要な量が揃わないため検定できない。

\begin{center}
\begin{tabularx}{\textwidth}{@{}llX@{}}
\toprule
必要な量 & 記号 & 状態 \\
\midrule
成果連動の強度 & $b$ & 返戻金の率と期間。形式が事業者ごとに異なる \\
結果のノイズ & $\sigma^2$ & 制御不能な変動のみを取り出す必要がある \\
努力の情報量 & 分子 & 代理変数を構成できていない \\
リスク回避度 & $r$ & 実験的手法を要する \\
\bottomrule
\end{tabularx}
\captionof{table}{$b^\star$ の検定に必要な量とその状態。}
\label{tab:bstar-requirements}
\end{center}

$\sigma^2$ は集計すれば足りるものではない。
必要なのは\textbf{事業者の努力では制御できない部分の分散}である。
職種別の離職率分散には事業者の選別や定着支援の巧拙が混入しており、
制御可能な変動と制御不能な変動が分離されていない。

そこで代わりに\textbf{価格側の反応}を見る。
可測性仮説が正しければ、成果が不確実な職種ほど
事業者は手数料を低く設定せざるをえないはずである。

\begin{quote}
\textbf{予測X}：職種別の手数料実績率は、その職種の離職率と負に相関する。

\textbf{対抗仮説}：手数料率は職種の年収水準や商慣行で決まり、離職率とは無関係である。
\end{quote}

\subsection{検定の結果}
\label{sec:predX-result}

人材サービス総合サイトの詳細検索により、
職種・手数料実績率・離職率を同時に取得した。
標本は東京都、許可番号13-ユ、四職種である。

\begin{center}
\begin{tabular}{@{}lrrrrrrl@{}}
\toprule
職種 & $n$ & 手数料率 & 中央値 & SD & 離職率 & $r$ & 慣行値の割合 \\
\midrule
看護師 & 82 & 23.0\% & 21.9 & 6.9 & 12.8\% & \textminus 0.171 & 40\% \\
医師 & 101 & 22.9\% & 20.0 & 6.2 & 2.3\% & +0.015 & 64\% \\
施設介護 & 67 & 22.2\% & 22.1 & 5.3 & 6.4\% & \textminus 0.103 & 45\% \\
保育士 & 57 & 25.6\% & 25.0 & 5.3 & 5.8\% & +0.025 & 54\% \\
\midrule
計 & 307 & 23.3\% & 22.2 & 6.1 & 6.7\% & \textminus 0.095 & --- \\
\bottomrule
\end{tabular}
\captionof{table}{職種別の手数料実績率と離職率、予測Xの検定結果（307社）。}
\label{tab:predX-result}
\end{center}
\begin{center}\footnotesize 慣行値の割合は手数料率が 20/25/30\% ちょうどである事業者の比率\end{center}

標本の範囲は注\ref{rem:predX-scope}に、
副次的な観測は注\ref{rem:scale-turnover}に記す。

\textbf{予測Xは支持されない。}
四職種すべてで相関は有意でなく（$|t|\leq 1.55$）、符号も一定しない。
307社をプールしても $r=-0.095$、$t=-1.67$ である。

規模を統制しても変わらない。
無期就職者数の対数を統制変数とした偏相関は $-0.084$（$t=-1.19$）、
さらに職種平均を除去した残差では $-0.055$（$t=-0.78$）である。

\subsection{職種間の非対称}
\label{sec:fee-rigidity}

より重要なのは職種間の対比である。

\begin{center}
\begin{tabular}{@{}lrr@{}}
\toprule
 & 範囲 & 比 \\
\midrule
手数料率（職種平均） & 22.2 -- 25.6\% & 1.16 倍 \\
離職率（職種平均） & 2.3 -- 12.8\% & \textbf{5.49 倍} \\
\bottomrule
\end{tabular}
\captionof{table}{手数料率と離職率、職種間のばらつきの比較。}
\label{tab:fee-turnover-asymmetry}
\end{center}

\textbf{離職率が5.5倍異なるにもかかわらず、手数料率は1.16倍しか異ならない。}
職種レベルの相関も $r=-0.102$ とほぼゼロである。

さらに手数料率は少数の値に集中する。
20\%、25\%、30\% のいずれかちょうどである事業者が全体の 40--64\% を占め、
15--35\% の範囲に 9割が収まる。
定着が最も安定している医師で集中度が最も高い（64\%）。

\begin{remark}[理論が想定する調整が作動していない]
\label{rem:no-adjustment}
例\ref{ex:bstar}の $b^\star=1/(1+r\sigma^2 k)$ は、
結果のノイズに応じて成果連動の強度を調整することを予測する。
本節の観測は、\textbf{その調整が価格に現れていない}ことを示す。

これは第\ref{sec:obstacle-taxonomy}節の四分類のいずれとも異なる。

\begin{center}
\begin{tabularx}{\textwidth}{@{}lX@{}}
\toprule
（i）測定の不在 & 該当しない。307社分が公表されている \\
（ii）仮説の不良設定 & 別途該当するが、本節の観測はそれとは独立 \\
（iii）識別不能 & 該当しない。交絡は統制済み \\
（iv）アクセス制約 & 該当しない。データに到達している \\
\bottomrule
\end{tabularx}
\captionof{table}{予測Xの観測結果と検証障害の四分類との対応（いずれにも該当しない）。}
\label{tab:predX-taxonomy-check}
\end{center}

説明変数が理論の想定する範囲で変動していない。
価格が慣行で決まるならば、$\sigma^2$ をどれだけ精密に測っても $b$ は反応しない。
\textbf{理論の精密化でもデータの追加でも解決しない}という点で、
これまでのどの障害とも性質が異なる。

同時にこれは記述的な発見でもある。
医療・介護・保育分野の人材紹介において、
手数料率は成果のリスクに応じて調整されておらず、
20--30\% の慣行値に集中している。
\end{remark}

\begin{remark}[本節の標本の限界]
\label{rem:predX-scope}
標本は\textbf{東京都、許可番号13-ユ、四職種}に限られる。全国を代表しない。

第一に、全国展開する事業者の本社が東京に集中するため、
大規模事業者が過剰に含まれる。無期就職1,000人超の事業者が複数ある。

第二に、他県労働局で許可を得た事業者の東京事業所が除外されている。

第三に、\textbf{詳細検索で職種を指定できるのは医療・介護・保育分野のみ}である。
情報通信や営業といった他分野では、同じ検定が構成できない。

第四に、手数料を金額で表示する事業者を除外している。
四職種で 307社中の約四分の一がこれにあたる。

なお標本が東京であることは、結論を弱める方向には働きにくい。
事業者密度が全国で最も高い地域であり、
\textbf{競争が最も激しい市場でも慣行値への集中が観測される}ことになる。
\end{remark}

\begin{remark}[副次的な観測：規模と離職率]
\label{rem:scale-turnover}
規模（無期就職者数の対数）と離職率の相関は職種で大きく異なる。

\begin{center}
\begin{tabular}{@{}lr@{}}
\toprule
職種 & $r$（規模, 離職率） \\
\midrule
看護師 & \textminus 0.009 \\
医師 & +0.191 \\
施設介護 & +0.431 \\
保育士 & +0.331 \\
\bottomrule
\end{tabular}
\captionof{table}{職種別の規模（無期就職者数の対数）と離職率の相関。}
\label{tab:scale-turnover-corr}
\end{center}

介護と保育で強い正の相関がある。大規模な事業者ほど離職率が高い。
両分野は適正認定制度の対象であり（第\ref{sec:regulated-regressor}節）、
大量紹介と定着の両立が難しい構造が示唆される。
本稿はこれを検証していない。
\end{remark}

\subsection{B1 の代理変数の候補：規模と離職率}
\label{sec:b1-proxy}

表\ref{tab:bstar-requirements}で、努力の情報量すなわちB1側の代理変数を
構成できていないと記した。注\ref{rem:scale-turnover}の観測がその候補になりうる。

\paragraph{機構。}
規模の拡大が一件あたりの投入を薄めるならば、
\textbf{事業者の努力が定着に効く職種でのみ離職率が上がる}。
努力が効かない職種では、投入が薄まっても離職率は変わらない。
したがって職種別の「規模と離職率の相関」は、B1の大きさの順序を与える。

観測された順序は次のとおりである（注\ref{rem:scale-turnover}）。

\begin{center}
\begin{tabular}{@{}lrl@{}}
\toprule
職種 & $r$（規模, 離職率） & B1 の含意 \\
\midrule
施設介護 & +0.431 & 努力が効く \\
保育士 & +0.331 & 同上 \\
医師 & +0.191 & 中間 \\
看護師 & \textminus 0.009 & 努力が効かない \\
\bottomrule
\end{tabular}
\captionof{table}{規模と離職率の相関を B1 の代理と読んだ場合の職種順序。}
\label{tab:b1-proxy-order}
\end{center}

\begin{remark}[機構は識別されない]
\label{rem:b1-proxy-assumption}
上の読みは\textbf{規模の拡大が一件あたりの投入を薄めるという仮定}に依存する。
容量制約が規模に比例して拡大するならば、投入は薄まらず相関は生じない。
本稿はこの仮定を検証していない。

また同じ符号は選別の緩みからも生じる。
大規模な事業者ほど求職者を多く流し、
一件あたりの適合の確認が粗くなるという経路である。
\textbf{投入の希薄化と選別の緩みは、この観測では区別できない。}
いずれもB1と同方向に働くため代理変数としては用いうるが、機構の主張はできない。

なお規模と紹介の質の関係には労働経済学に文献があるが、本稿は渉猟していない。
ここで述べるのは新規性ではなく、既にある観測がB1の代理として使えるかという点である。
\end{remark}

\paragraph{予測の登録。}
第\ref{ch:next}章の手続きに従い、検定の前に次を固定する。

\begin{quote}
\textbf{予測Y}：職種別の「規模と離職率の相関」は、
他都道府県の標本においても表\ref{tab:b1-proxy-order}の順序を保つ。
すなわち施設介護と保育士で正、看護師でゼロ近傍である。

\textbf{対抗仮説}：相関は東京都に固有の事情、
とくに全国展開する事業者の本社が集中していることによる標本の産物であり、
他県では現れない。
\end{quote}

\begin{quote}
\textbf{予測Z}：規模と離職率の相関が強い職種ほど、返戻金制度が強い。
B1が大きい職種では $\pi$ を $y$ に載せる価値が高いためである。

\textbf{対抗仮説}：返戻金の強度は離職率の水準に対応する。
すなわちB2（リスク側）が支配し、離職率の高い職種ほど返戻金は弱い。
\end{quote}

予測Zが成立すれば、第\ref{sec:predB-illposed}節で符号が一意に定まらないとした
予測Bについて、B1とB2を分離する経路が得られる。
予測Yはその前提であり、相関が標本の産物でないことを確認する。

\paragraph{事前に確定した事項。}

\begin{center}
\begin{tabularx}{\textwidth}{@{}lX@{}}
\toprule
項目 & 内容 \\
\midrule
母集団 & 人材サービス総合サイトに事業報告を掲載する事業所。都道府県を指定し、許可番号の接頭辞で識別する \\
職種 & 詳細検索で指定できる医療・介護・保育分野。表\ref{tab:predX-result}の四職種を最低限含める \\
規模の定義 & 無期雇用就職者数の対数。第\ref{sec:predX-result}節と同一 \\
下限 & 無期雇用就職者数が 0 の事業所は離職率が定義されないため除く \\
除外 & 手数料を金額で表示する事業所は予測Zでのみ除く。予測Yでは用いる \\
同定 & 検索結果からの取得に限る。直接指定は制限されている（第\ref{sec:pilot}節） \\
\bottomrule
\end{tabularx}
\captionof{table}{予測YとZについて事前に確定した事項。}
\label{tab:predYZ-registration}
\end{center}

\textbf{本稿はこれを検定していない。} 第\ref{sec:predX-result}節の標本は東京都に限られ、
他都道府県の取得を要する。取得手続きは同一であるため、
第\ref{sec:search-route}節の経路をそのまま適用できる。

\begin{remark}[同一標本で検定してはならない]
\label{rem:no-same-sample-test}
表\ref{tab:b1-proxy-order}の相関は307社の標本において\emph{事後に}観測されたものである。
同じ標本で予測Yを検定すれば、
第\ref{ch:method}章の事後的物語化と同型の誤りになる。
予測Yが要求するのは標本外での再現であり、
これが満たされない限り相関は照合にとどまる（第\ref{sec:comparison-vs-test}節）。
\end{remark}

\subsection{予測Yの検定}
\label{sec:predY-result}

人材サービス総合サイトの詳細検索を、都道府県を全国として四職種について取得した。
第\ref{sec:predX-result}節の標本は東京都に限られていたため、
\textbf{東京都を除いた部分が標本外にあたる}。

\begin{center}
\begin{tabular}{@{}lrrr@{}}
\toprule
職種 & 検索結果 & 抽出行数 & 名寄せ後 \\
\midrule
施設介護の職業 & 1,075 & 1,075 & 464 \\
保育士 & 552 & 552 & 200 \\
医師 & 479 & 479 & 269 \\
看護師、准看護師 & 1,322 & 1,322 & 469 \\
\bottomrule
\end{tabular}
\captionof{table}{予測Yに用いた標本。単位は事業所行、名寄せ後は企業。}
\label{tab:predY-sample}
\end{center}

抽出行数は検索結果と一致する。許可番号の記号は
ユ（有料）3,221 件、ム（無料職業紹介所）200 件、特 7 件の三種である。

\begin{remark}[行は事業所単位、値は企業単位である]
\label{rem:jinzai-dedup}
検索結果の各行は事業所を表すが、
\textbf{就職者数・離職率・手数料は許可番号すなわち企業単位の値であり、
その企業の全事業所に同じ値が複製されている}。

医師では株式会社パソナ一社が479行中134行を占める。
名寄せせずに相関を取ると同社が134倍の重みを持ち、
無期就職者数の中央値は14人から6,430人に、相関の符号は正から負に変わる。
\textbf{許可番号による名寄せは前処理ではなく、結論を決める手続きである。}
\end{remark}

\paragraph{基準値の再現。}
表\ref{tab:b1-proxy-order}の相関を、今回の標本の東京都部分で再現する。

\begin{center}
\begin{tabular}{@{}lrrrr@{}}
\toprule
職種 & 表\ref{tab:b1-proxy-order} & $n$ & 今回の東京 & $n$ \\
\midrule
施設介護 & +0.431 & 46 & +0.423 & 47 \\
保育士 & +0.331 & 38 & +0.349 & 40 \\
医師 & +0.191 & 49 & +0.186 & 50 \\
看護師 & \textminus 0.009 & 69 & \textminus 0.009 & 70 \\
\bottomrule
\end{tabular}
\captionof{table}{基準値の再現。標本数は四職種とも2以内で一致する。}
\label{tab:predY-baseline}
\end{center}

\textbf{再現されたとみなす。}

\paragraph{結果。}
東京都を除いた標本で同じ量を算出する。

\begin{center}
\begin{tabular}{@{}lrrrrrr@{}}
\toprule
 & \multicolumn{3}{c}{登録どおり（金額表示を含む）} & \multicolumn{3}{c}{金額表示を除く} \\
\cmidrule(lr){2-4}\cmidrule(lr){5-7}
職種 & $n$ & $r$ & $t$ & $n$ & $r$ & $t$ \\
\midrule
施設介護 & 176 & +0.246 & \textbf{+3.34} & 122 & +0.221 & \textbf{+2.48} \\
保育士 & 59 & +0.163 & +1.25 & 42 & +0.645 & \textbf{+5.34} \\
医師 & 54 & +0.159 & +1.16 & 43 & +0.126 & +0.81 \\
看護師 & 150 & +0.067 & +0.81 & 116 & +0.070 & +0.75 \\
\midrule
順位 & \multicolumn{3}{l}{施設介護 $>$ 保育士 $>$ 医師 $>$ 看護師} & \multicolumn{3}{l}{保育士 $>$ 施設介護 $>$ 医師 $>$ 看護師} \\
\bottomrule
\end{tabular}
\captionof{table}{東京都を除いた標本における規模と離職率の相関。}
\label{tab:predY-result}
\end{center}

\paragraph{判定。}
事実を分けて述べる。

\begin{enumerate}[label=(\arabic*)]
\item \textbf{登録した規則の下では、登録した内容がすべて満たされる。}
順位は表\ref{tab:b1-proxy-order}と完全に一致し、
施設介護と保育士は正、看護師はゼロ近傍である。

\item \textbf{有意なのは施設介護のみである。}
他の三職種は $|t|<1.3$ であり、順位の差を相関の差として読むことはできない。

\item \textbf{順位の一致は規則に依存する。}
金額表示を除くと保育士が施設介護を上回り、上位二つが入れ替わる。

\item \textbf{定性的な区分は規則によらない。}
施設介護と保育士が正、看護師がゼロ近傍である点は両規則で成立する。
\end{enumerate}

\textbf{判定は部分的支持とする。} (1)により支持と読むこともできるが、
注\ref{rem:predY-registration-error}のとおり\textbf{登録した規則が
基準値の構成と一致していなかった}。
規則が事前に正しく固定されていない以上、
(3)で有利に出る側を事前に選んだとは言えない。
\textbf{規則によらずに成立するのは(4)の定性的な区分のみである。}

\begin{remark}[登録した除外規則が基準値の構成と一致していなかった]
\label{rem:predY-registration-error}
表\ref{tab:predYZ-registration}は、金額表示の事業所を予測Yでは用いると定めた。
しかし表\ref{tab:b1-proxy-order}の基準値は
金額表示を除いた標本について算出されている（注\ref{rem:predX-scope}）。
\textbf{登録した規則は、比較の対象である基準値を生む構成ではなかった。}

両方を報告し、判定を規則によらない部分に限ったのはこのためである。
\textbf{登録の時点で基準値の構成を確認していれば生じなかった不備である。}
\end{remark}

\begin{remark}[標本の限界]
\label{rem:predY-limits}
二点を記す。

第一に、表示される離職率は同じ行の就職者数と離職者数から計算できない。
医師では一致するのが479行中115行にとどまり、両者は期間が異なる。
離職率はサイト側の集計値として用いている。
これは第\ref{sec:predX-result}節と同じ構成であり比較には支障しないが、
\textbf{規模と離職率が同一期間の量ではないことは記録しておく}。

第二に、東京都を除いた標本では規模の分布が異なる。
無期就職者数の中央値は東京都で23人、それ以外で5人である。
\end{remark}

\section{判定}
\label{sec:family-assignment}

\begin{quote}
族2 の主要障害は（iv）アクセス制約であり、必要な粒度のデータに到達できない。
族5 の主要障害は、$b^\star$ の式については（ii）仮説の不良設定、
価格側の検定については（vi）調整の不作動である。
\end{quote}

族5 では、経路確認の結果データ源は良好であった（第\ref{sec:jinzai-structure}節）。
実際、価格側からの検定は307社について実施できた（第\ref{sec:predX-result}節）。
それでも検証が進まないのは、
$b^\star$ の式については命題が符号を一意に定めないこと、
価格側については説明変数が理論の想定する範囲で変動しないことによる。

この判別には実務上の含意がある。
（iii）（iv）（v）が障害であれば労力の投入は前進を生むが、
（i）（ii）（vi）が障害であれば同じ労力は何も生まない。
\textbf{どちらであるかを先に判定しないと、労力の配分を誤る。}

\section{取得の経路}
\label{sec:pilot}

族5 のデータ取得は二つの経路を経た。
第一の経路は事業所ごとの詳細情報を個別に開く方法であり、これは中止した。
許可番号を含むURLを直接指定する取得では、
\textbf{解決に失敗した番号に対してエラーではなく直前の結果が返ることがあった}。
自動的な取得を制限する措置と考えられ、本稿はこの挙動を分析の対象としない。

第二の経路は詳細検索の条件指定であり、こちらが機能した。以下はその記録である。

\subsection{第二の経路：詳細検索}
\label{sec:search-route}

同サイトの詳細検索は、職種・手数料実績率・離職率を検索条件に指定でき、
指定した場合はそれらが検索結果の一覧に表示される。
\textbf{個別の詳細情報を開く必要がない。}

この経路により、第\ref{sec:predX-result}節の307社分のデータを取得した。
第一の経路で遭遇した同定失敗は、
検索結果からの遷移では生じない。

ただし制約がある。
検索条件に指定できる職種は医療・介護・保育分野の9種に限られ、
情報通信や営業といった他分野では同じ取得ができない。

\subsection{当初の三経路の帰結}

当初の調査計画（第\ref{part:method}部）で挙げた三経路について、現時点の状態を記録する。

\paragraph{掲載率の測定。}
四職種の検索結果では、手数料を率で表示する事業者が約四分の三、
金額で表示する事業者が約四分の一であった。
掲載自体はほぼ全事業者で行われている。
義務化初年度としては入力が進んでいる。

\paragraph{分野別離職率の記述。}
\textbf{達成された。} 第\ref{sec:predX-result}節のとおり、
四職種の離職率が事業者単位で得られた。
職種平均は医師 2.3\%、保育士 5.8\%、施設介護 6.4\%、看護師 12.8\% であり、
\textbf{5.5倍の開きがある}。

\paragraph{規制の拘束度の測定。}
未達である。返戻金の掲載形式が事業者ごとに不統一であり、
一件ずつ人が読む作業になる。
ただし手数料率については、認定事業者の標準偏差が 2.2 と
非認定の 7.0 より小さいという観測が得られた（第\ref{sec:predX-result}節）。
認定条件に手数料率の規定はないため、これは間接的な効果である。

\subsection{記録の粒度の不整合}

第\ref{sec:unit-mismatch}節で指摘した企業と事業所の粒度の不整合は、
第\ref{sec:predY-result}節の全国の標本において明確な形で現れる
（注\ref{rem:jinzai-dedup}）。検索結果の行は事業所を表すが、
就職者数・離職率・手数料は企業単位の値であり、その企業の全事業所に複製されている。

\section{四分類に収まらない障害}
\label{sec:sixth-obstacle}

本章の検討により、第\ref{sec:obstacle-taxonomy}節の四分類に収まらない障害が二つ現れた。
いずれも本稿の調査で実際に遭遇したものである。

\begin{center}
\begin{tabularx}{\textwidth}{@{}lXl@{}}
\toprule
種別 & 内容 & 現れた箇所 \\
\midrule
（v）集計の選択 & 調査はされるが公表される集計表に含まれない & 第\ref{sec:not-tabulated}節 \\
（vi）調整の不作動 & データは十分にあるが、説明変数が理論の想定する範囲で変動しない & 注\ref{rem:no-adjustment} \\
\bottomrule
\end{tabularx}
\captionof{table}{四分類に収まらない障害（v・vi）とその現れた箇所。}
\label{tab:sixth-obstacle-summary}
\end{center}

両者は対処が異なる。

（v）は集計要望や匿名データの利用申請という経路がありうる。
標本には含まれている可能性が高いため、公表の範囲が変われば解決する。

（vi）は\textbf{どの経路でも解決しない}。
理論を精密化しても、データを増やしても、
説明変数が動いていない以上は検定が成立しない。
これは理論の欠陥でも測定の欠陥でもなく、
\textbf{理論が想定する調整機構が現実に作動していないという事実}である。

\begin{remark}[（vi）は否定的な結果ではない]
理論の予測が現れないことは、それ自体が観測である。
注\ref{rem:no-adjustment}の場合、
「手数料率が成果のリスクに応じて調整されていない」という記述が得られた。

第\ref{sec:layers}章のレイヤー構造でいえば、
価格が慣行で決まることは制度層に属する現象である。
最適化を前提とする理論が、慣行の前で作動しないという構図になる。
\end{remark}

\subsection{（v）の実例：個人事業主の信用ポジション}
\label{sec:not-tabulated}

個人事業主の $\kappa$ を測ろうとした際、二つの統計に当たった。

\begin{center}
\begin{tabularx}{\textwidth}{@{}lXl@{}}
\toprule
統計 & 状態 & 判定 \\
\midrule
個人企業経済調査 & 資産・負債を調査していたが2019年度に廃止 & 測定の\emph{喪失} \\
中小企業実態基本調査 & 個人事業者用の調査票は存在するが、資産・負債の集計表は法人企業のみ & 調査されるが\emph{集計されない} \\
\bottomrule
\end{tabularx}
\captionof{table}{個人事業主の $\kappa$ に関わる二つの統計の状態。}
\label{tab:individual-kappa-stats}
\end{center}

前者は（i）測定の不在だが、\textbf{かつては測定されていた}点で通常の事例と異なる。
2018年度以前のデータは存在するため、過去については分析可能である。

後者はより特異である。調査票が存在する以上、
回答は収集されている可能性が高い。しかし公表される集計表に含まれない。
これは（i）でも（iv）でもなく、\textbf{集計の選択}という第五の障害である。

\begin{remark}[集計の選択が観測を規定する]
\label{rem:tabulation-choice}
第\ref{part:method}部で扱った選択バイアスは、
標本の選択と変数の測定に関するものであった。
本節が示すのは\textbf{集計の段階での選択}である。

どの断面を公表するかは統計作成者が決める。
法人企業について資産・負債を集計し個人企業について集計しないという判断は、
中小企業施策の必要性に基づくものであろうが、
結果として\textbf{個人事業主の信用ポジションは観測されない}。

第\ref{sec:frame}節で述べた「フレームの網羅性は標本規模を保証しない」に加え、
\textbf{標本に含まれることは集計されることを保証しない}。
\end{remark}



\section{現時点の到達点と断念}

\subsection{断念した対象}

\begin{enumerate}[label=(\arabic*)]
\item \textbf{本稿のデータによる}族2 の $\phi_{\mathrm{cog}}$ の測定。
命題\ref{prop:phicog-requirements}の (2)(3) を欠くため、
個票の利用記録と外生ショックを伴うデータに到達しない限り推定できない。
なお同種の量は他者により測定されている（第\ref{sec:phicog-measurable}節）。
断念するのは\emph{公表統計による測定}であって、測定そのものではない。

\item 族5 の $b^\star$ の式の検定。
$b$ を連続量として得ること、および $\sigma^2$ から
制御不能な部分のみを取り出すことができない（（ii）仮説の不良設定）。

\item 価格側からの検定の続行。
第\ref{sec:predX-result}節により無相関が確認された。
説明変数が慣行値に集中して変動しないため、
標本を増やしても結果は変わらない（（vi）調整の不作動）。

\item 事業所詳細情報の個別取得による大規模収集。
第\ref{sec:pilot}節のとおり直接指定による取得は制限されている。
詳細検索による取得はこの制約を受けない（第\ref{sec:search-route}節）。
\end{enumerate}

\subsection{条件付きで可能な対象}

\begin{enumerate}[label=(\arabic*)]
\setcounter{enumi}{4}
\item 他都道府県への拡張。
第\ref{sec:predX-result}節の標本は東京都に限られる（注\ref{rem:predX-scope}）。
詳細検索は都道府県を指定するため、他県での再現は同じ手続きで可能である。

\item 規模と離職率の関係の検証。
第\ref{sec:b1-proxy}節で予測YとZを登録した。
検定には他都道府県の標本を要する。

\item 返戻金の率と期間の収集。
事業者ごとに掲載形式が異なるため、一件ずつ人が読む作業になる。
認定事業者60--70社に限れば実行可能な規模である。
\end{enumerate}

\subsection{族2と族5の到達点}

\begin{center}
\begin{tabularx}{\textwidth}{@{}lXX@{}}
\toprule
 & 族2 & 族5 \\
\midrule
検証できたこと & $\kappa$ の水準差が二桁。対抗仮説を部分的に分離
 & 手数料率と離職率が無相関。慣行値への集中 \\
検証できなかったこと & $\phi_{\mathrm{cog}}$（本稿のデータでは） & $b^\star$ の式 \\
主要障害 & （iv）アクセス制約 & （ii）不良設定と（vi）調整の不作動 \\
副次的障害 & 個社データの守秘義務 & 記録の粒度、母集団の小ささ、同定失敗 \\
再開の条件 & 個票データへの到達 & なし。（vi）は労力で解決しない \\
\bottomrule
\end{tabularx}
\captionof{table}{族2と族5、それぞれの到達点のまとめ。}
\label{tab:family2-5-summary}
\end{center}

族2 については、当初の予測のうち水準に関する部分が支持された。

族5 については、$b^\star$ の式の検定には至らなかったが、
\textbf{価格側の検定により記述的な結果が得られた}。
医療・介護・保育分野の人材紹介において、
職種平均の離職率が5.5倍異なるにもかかわらず手数料率は1.16倍しか異ならず、
20--30\% の慣行値に集中している。

第\ref{sec:predB-illposed}節で仮説の構造的欠陥を特定したことに加え、
第\ref{sec:predX-result}節で\textbf{理論が想定する調整が現実に作動していない}ことを
示せたのは、データの追加では到達しえなかった結果である。

\section{枠組みへの追加改訂}

第\ref{part:results}部の改訂に加えて一点を記録する。

本章の四分類は、検証が完結しなかった原因の事後の分類であると同時に、
\textbf{調査に着手する前の判定にも使える}。
第\ref{ch:design}章の段階3において、
当該の量が必要な粒度で取得できるかを判定する際の基準がこれにあたる。
本章の位置づけは、実証の記録であると同時に設計の道具でもある。

第\ref{ch:platformfee}章の調査可能性の判定は、この使い方の実例である。
